\documentclass{ximera}
\input{../preamble.tex}

\title{Orthogonal Matrices and Symmetric Matrices} \license{CC BY-NC-SA 4.0}

\begin{document}

\begin{abstract}

\end{abstract}
\maketitle

\begin{onlineOnly}
\section*{Orthogonal Matrices and Symmetric Matrices}
\end{onlineOnly}

Recall that an $n \times n$ matrix $A$ is diagonalizable if and only if it has $n$ linearly independent eigenvectors.  Moreover, the matrix $P$ with these eigenvectors as columns is a diagonalizing matrix for $A$, that is
\begin{equation*}
P^{-1}AP =D
\end{equation*}
where $D$ is a diagonal matrix with diagonal entries consisting of the eigenvalues of $A$.
% As we have seen, the nice bases of $\RR^n$ are the orthogonal ones, so a natural question is: which $n \times n$ matrices have $n$ orthogonal eigenvectors, so that columns of $P$ form an orthogonal basis for $\RR^n$? These turn out to be precisely the \dfn{symmetric matrices} (matrices for which $A=A^T$), and this is the main result of this section.

\section*{Orthogonal Matrices}
A collection of non-zero, pairwise orthogonal vectors in $\RR^n$ is called an \dfn{orthogonal} set of vectors. An orthogonal set of vectors is called \dfn{orthonormal} if $\norm{\vec{q}} = 1$ for each vector $\vec{q}$ in the set. A set of orthogonal vectors $\{\vec{v}_{1}, \vec{v}_{2}, \dots, \vec{v}_{k}\}$ can be ``\textit{normalized}'', i.e. converted into an orthonormal set $\left\{ \frac{1}{\norm{\vec{v}_{1}}}\vec{v}_{1}, \frac{1}{\norm{\vec{v}_{2}}}\vec{v}_{2}, \dots, \frac{1}{\norm{\vec{v}_{k}}}\vec{v}_{k} \right\}$. In particular, if a matrix $A$ has $n$ orthogonal eigenvectors, they can (by normalizing) be taken to be orthonormal. The corresponding diagonalizing matrix (we will use $Q$ instead of $P$) has orthonormal columns, and such matrices are very easy to invert.


\begin{theorem}\label{th:orthogonal_matrices}
The following conditions are equivalent for an $n \times n$ matrix $Q$.

\begin{enumerate}
\item\label{th:orthogonal_matrices_a} $Q$ is invertible and $Q^{-1} = Q^{T}$.

\item\label{th:orthogonal_matrices_b} The rows of $Q$ are orthonormal.

\item\label{th:orthogonal_matrices_c} The columns of $Q$ are orthonormal.

\end{enumerate}
\end{theorem}

\begin{proof}
First note that condition \ref{th:orthogonal_matrices_a} is equivalent to $Q^{T}Q = I$. Let $\vec{q}_{1}, \vec{q}_{2}, \dots, \vec{q}_{n}$ denote the columns of $Q$. Then $\vec{q}_{i}^{T}$ is the $i$th row of $Q^{T}$, so the $(i, j)$-entry of $Q^{T}Q$ is $\vec{q}_{i} \dotp \vec{q}_{j}$. Thus $Q^{T}Q = I$ means that $\vec{q}_{i} \dotp \vec{q}_{j} = 0$ if $i \neq j$ and $\vec{q}_{i} \dotp \vec{q}_{j} = 1$ if $i = j$. Hence condition \ref{th:orthogonal_matrices_a} is equivalent to \ref{th:orthogonal_matrices_c}. The proof of the equivalence of \ref{th:orthogonal_matrices_a} and \ref{th:orthogonal_matrices_b} is similar.
\end{proof}

\begin{definition}[Orthogonal Matrices]\label{def:orthogonal matrices}
An $n \times n$ matrix $Q$ is called an \dfn{orthogonal matrix} if its columns are orthonormal vectors.  (i.e. The matrix satisfies one (and hence all) of the conditions in Theorem~\ref{th:orthogonal_matrices}.)
\end{definition}

\begin{remark}\label{rem:orthVsOrthnormMat}
The above definition calls a square matrix with orthonormal columns an \dfn{orthogonal matrix}.  It is true that \dfn{orthonormal matrix} might be a better name. But \dfn{orthogonal matrix} is standard.
\end{remark}


\begin{example}\label{ex:rotation_ortho}
The rotation matrix
$\begin{bmatrix}
\cos\theta & -\sin\theta \\
\sin\theta & \cos\theta
\end{bmatrix}$ is orthogonal for any angle $\theta$.
\begin{explanation}
See Practice Problem \ref{prob:rotation_ortho}.
\end{explanation}
\end{example}




%\begin{exploration}\label{exp:make_orthogonal}
%\begin{octave} 
%A=[2 1 1; -1 1 1; 0 -1 1]
%\end{octave}
%\begin{enumerate}
 %   \item Check that matrix $A$ in the Octave window has rows that are orthogonal.
  %  \begin{hint}
  %  A(1,:)*transpose(A(2,:)) 
    
  %  A(2,:)*transpose(A(3,:)) 
    
  %  A(1,:)*transpose(A(3,:)) 
    
  %  \end{hint}
 %   \item Check that matrix $A$ in the Octave window has columns that are NOT orthogonal.
 %   \begin{hint}
 %   transpose(A(:,1))*A(:,2) 
    
 %   etc.
    
 %   \end{hint}
 %   \item Check that matrix $A$ in the Octave window has rows that are NOT orthonormal.
 %   \begin{hint}
 %   A(1,:)*transpose(A(1,:)) 
    
 %   A(2,:)*transpose(A(2,:)) 
    
 %   A(3,:)*transpose(A(3,:)) 
    
 %   \end{hint}
%    \item Create a matrix $Q$ by normalizing each of the rows of $A$.
%    \begin{hint}
%    q1=A(1,:)/norm(A(1,:)); 
    
%    q2=A(2,:)/norm(A(2,:));
    
%    q3=A(3,:)/norm(A(3,:));
    
%    Q = [q1;q2;q3]
    
%    \end{hint}
%    \item Check that $Q$ is an orthogonal matrix.
%    \begin{hint}
%    You should get $Q = \begin{bmatrix}
%\frac{2}{\sqrt{6}} & \frac{1}{\sqrt{6}} & %\frac{1}{\sqrt{6}} \\
%\frac{-1}{\sqrt{3}} & \frac{1}{\sqrt{3}} & %\frac{1}{\sqrt{3}} \\
%0 & \frac{-1}{\sqrt{2}} & \frac{1}{\sqrt{2}}
%\end{bmatrix}$, and one can check that this is orthogonal in a number of ways.
 %   \end{hint}
%\end{enumerate}

%\end{exploration}

\begin{exploration}\label{exp:make_orthogonal}
Let $A=\begin{bmatrix}
    2&1&1\\-1&1&1\\0&-1&1
\end{bmatrix}$
\begin{enumerate}
    \item Check that matrix $A$ has rows that are orthogonal.
    \item Check that matrix $A$ has columns that are NOT orthogonal.
    \item Check that matrix $A$ has rows that are NOT orthonormal.
    \item Create a matrix $Q$ by normalizing each of the rows of $A$.
    \item Check that columns and rows of $Q$ are orthonormal.
\end{enumerate}


You should get $Q = \begin{bmatrix}
\frac{2}{\sqrt{6}} & \frac{1}{\sqrt{6}} & \frac{1}{\sqrt{6}} \\
\frac{-1}{\sqrt{3}} & \frac{1}{\sqrt{3}} & \frac{1}{\sqrt{3}} \\
0 & \frac{-1}{\sqrt{2}} & \frac{1}{\sqrt{2}}
\end{bmatrix}$.

\end{exploration}

\begin{definition}\label{def:orthDiag}
An $n \times n$ matrix $A$ is said to be \dfn{orthogonally diagonalizable} if an orthogonal matrix $Q$ can be found such that  $Q^{-1}AQ = Q^{T}AQ$ is diagonal.
\end{definition}

We have learned earlier that when we diagonalize a matrix $A$, we write $P^{-1}AP=D$ for some matrix $P$ where $D$ is diagonal, and the diagonal entries are the eigenvalues of $A$.  We have also learned that the columns of the matrix $P$ are the corresponding eigenvectors of $A$.  So when a matrix is orthogonally diagonalizable, we are able to accomplish the diagonalization using a matrix $Q$ consisting of $n$ eigenvectors that form an orthonormal basis for $\RR^n$.  In the following section we will learn that the matrices that have this property are precisely the symmetric matrices.

\section*{Symmetric Matrices}

A \dfn{symmetric matrix} is a matrix which is equal to its transpose.  

When we began our study of eigenvalues and eigenvectors, we saw examples of matrices with entries that were real numbers with eigenvalues that were complex numbers.  It can be shown that symmetric matrices only have real eigenvalues.  

\begin{fact}
    Symmetric matrices only have real eigenvalues.
\end{fact}

\begin{lemma}\label{th:dotpSymmetric}
If A is an $n \times n$ symmetric matrix, then
\begin{equation*}
(A\vec{x}) \dotp \vec{y} = \vec{x} \dotp (A\vec{y})
\end{equation*}
for all columns $\vec{x}$ and $\vec{y}$ in $\RR^n$.
\end{lemma}

\begin{proof}
Recall that $\vec{x} \dotp \vec{y} = \vec{x}^{T} \vec{y}$ for all columns $\vec{x}$ and $\vec{y}$. Because $A^{T} = A$, we get
\begin{equation*}
(A\vec{x}) \dotp \vec{y} = (A\vec{x})^T\vec{y} = \vec{x}^TA^T\vec{y} =  \vec{x}^TA\vec{y} = \vec{x} \dotp (A\vec{y})
\end{equation*}
\end{proof}

\begin{remark}
The converse of the above lemma also holds (see Practice Problem \ref{ex:8_2_15}).
\end{remark}

\begin{theorem}\label{th:symmetric_has_ortho_ev}
If $A$ is a symmetric matrix, then eigenvectors of $A$ corresponding to distinct eigenvalues are orthogonal.
\end{theorem}

\begin{proof}
Let $A\vec{x} = \lambda \vec{x}$ and $A\vec{y} = \mu \vec{y}$, where $\lambda \neq \mu$. We compute
\begin{equation*}
\lambda(\vec{x} \dotp \vec{y}) = (\lambda\vec{x}) \dotp \vec{y} = (A\vec{x}) \dotp \vec{y} = \vec{x} \dotp (A\vec{y}) = \vec{x} \dotp (\mu\vec{y}) = \mu(\vec{x} \dotp \vec{y})
\end{equation*}
Hence $(\lambda - \mu)(\vec{x} \dotp \vec{y}) = 0$, and so $\vec{x} \dotp \vec{y} = 0$ because $\lambda \neq \mu$.
\end{proof}

\begin{example}\label{ex:DiagonalizeSymmetricMatrix}
Find an orthogonal matrix $Q$ such that $Q^{-1}AQ$ is diagonal, where $A = \begin{bmatrix}
1 & 0 & -1 \\
0 & 1 & 2 \\
-1 & 2 & 5
\end{bmatrix}$.

\begin{explanation}
 The characteristic equation of $A$ is (adding twice row 1 to row 2):
\begin{equation*}
\det \begin{bmatrix}
1-\lambda & 0 & 1 \\
0 & 1-\lambda & -2 \\
1 & -2 & 5-\lambda
\end{bmatrix} = \lambda(\lambda - 1)(\lambda - 6)=0
\end{equation*}
Thus the eigenvalues are $\lambda = 0$, $1$, and $6$, and corresponding eigenvectors are
\begin{equation*}
\vec{x}_{1} = \begin{bmatrix}
1 \\
-2 \\
1
\end{bmatrix} \;
\vec{x}_{2} = \begin{bmatrix}
2 \\
1 \\
0
\end{bmatrix} \;
\vec{x}_{3} = \begin{bmatrix}
-1 \\
2 \\
5
\end{bmatrix}
\end{equation*}
respectively. Moreover,these eigenvectors are \textit{orthogonal}. We have $\norm{\vec{x}_{1}}^{2} = 6$, $\norm{\vec{x}_{2}}^{2} = 5$, and $\norm{\vec{x}_{3}}^{2} = 30$, so
\begin{equation*}
P = \begin{bmatrix}
| & | &  | \\
\frac{1}{\sqrt{6}}\vec{x}_{1} & \frac{1}{\sqrt{5}}\vec{x}_{2} & \frac{1}{\sqrt{30}}\vec{x}_{3} \\
| & | &  | 
\end{bmatrix} = \frac{1}{\sqrt{30}} \begin{bmatrix}
\sqrt{5} & 2\sqrt{6} & -1 \\
-2\sqrt{5} & \sqrt{6} & 2 \\
\sqrt{5} & 0 & 5
\end{bmatrix}
\end{equation*}
is an orthogonal matrix. Thus $P^{-1} = P^{T}$ and
\begin{equation*}
P^TAP = \begin{bmatrix}
0 & 0 & 0 \\
0 & 1 & 0 \\
0 & 0 & 6
\end{bmatrix}.
\end{equation*}

\end{explanation}
\end{example}


\begin{example}\label{exa:ortho_diag_symm}
Orthogonally diagonalize the symmetric matrix $A = \begin{bmatrix}
8 & -2 & 2 \\
-2 & 5 & 4 \\
2 & 4 & 5
\end{bmatrix}$.


\begin{explanation}
  The characteristic equation is
\begin{equation*}
\det  \begin{bmatrix}
8-\lambda & 2 & -2 \\
2 & 5-\lambda & -4 \\
-2 & -4 & 5-\lambda
\end{bmatrix} = \lambda(\lambda-9)^2=0
\end{equation*}
Hence the distinct eigenvalues are $0$ and $9$ are of algebraic multiplicity $1$ and $2$, respectively.  The geometric multiplicities must be the same, for $A$ is diagonalizable, being symmetric. It follows that $\mbox{dim}(\mathcal{S}_0) = 1$ and $\mbox{dim}(\mathcal{S}_9) = 2$.  Gaussian elimination gives
\begin{equation*}
\mathcal{S}_{0}(A) = \mbox{span}\{\vec{x}_{1}\}, \quad \vec{x}_{1} = \begin{bmatrix}
1 \\
2 \\
-2
\end{bmatrix}, \quad \mbox{ and } \quad \mathcal{S}_{9}(A) = \mbox{span} \left\lbrace \begin{bmatrix}
	-2 \\
	1 \\
	0
	\end{bmatrix}, \begin{bmatrix}
2 \\
0 \\
1
\end{bmatrix} \right\rbrace
\end{equation*}
The eigenvectors in $\mathcal{S}_{9}$ are both orthogonal to $\vec{x}_{1}$ as Theorem~\ref{th:symmetric_has_ortho_ev} guarantees, but not to each other. However, an orthogonal basis can be found using projections. (See \href{https://ximera.osu.edu/corela/LinearAlgebraInteractiveIntro/RTH-0015/main}{Gram-Schmidt Orthogonalization})
\begin{equation*}
\{\vec{f}_{2}, \vec{f}_{3}\} \mbox{ of } \mathcal{S}_{9}(A) \quad \mbox{ where } \quad \vec{f}_{2} = \begin{bmatrix}
-2 \\
1 \\
0
\end{bmatrix} \mbox{ and }  \vec{f}_{3} = \begin{bmatrix}
2 \\
4 \\
5
\end{bmatrix}
\end{equation*}
Normalizing gives orthonormal vectors $\{\frac{1}{3}\vec{x}_{1}, \frac{1}{\sqrt{5}}\vec{f}_{2}, \frac{1}{3\sqrt{5}}\vec{f}_{3}\}$, so
\begin{equation*}
Q = \begin{bmatrix}
| & | &  | \\
\frac{1}{3}\vec{x}_{1} & \frac{1}{\sqrt{5}}\vec{f}_{2} & \frac{1}{3\sqrt{5}}\vec{f}_{3} \\
| & | &  | 
\end{bmatrix} = \frac{1}{3\sqrt{5}}\begin{bmatrix}
\sqrt{5} & -6 & 2 \\
2\sqrt{5} & 3 & 4 \\
-2\sqrt{5} & 0 & 5
\end{bmatrix}
\end{equation*}
is an orthogonal matrix such that $Q^{-1}AQ$ is diagonal.


It is worth noting that other, more convenient, diagonalizing matrices $Q$ exist. For example, $\vec{y}_{2} = \begin{bmatrix}
2 \\
1 \\
2
\end{bmatrix}$ and $\vec{y}_{3} = \begin{bmatrix}
-2 \\
2 \\
1
\end{bmatrix}$
 lie in $\mathcal{S}_{9}(A)$ and they are orthogonal. Moreover, they both have norm $3$ (as does $\vec{x}_{1}$), so
\begin{equation*}
\hat{Q} = \begin{bmatrix}
| & | &  | \\
\frac{1}{3}\vec{x}_{1} & \frac{1}{3}\vec{y}_{2} & \frac{1}{3}\vec{y}_{3} \\
| & | &  |
\end{bmatrix} = \frac{1}{3}\begin{bmatrix}
1 & 2 & -2 \\
2 & 1 & 2 \\
-2 & 2 & 1
\end{bmatrix}
\end{equation*}
is a nicer orthogonal matrix with the property that $\hat{Q}^{-1}A\hat{Q}$ is diagonal.
\end{explanation}
\end{example}

\section*{Practice Problems}

\begin{problem}\label{prob:make_ortho_matrix3}
Normalize the columns to make the following matrix orthogonal.
$A = \begin{bmatrix}
1 & -4 \\
2 & 2
\end{bmatrix}$

$$\frac{1}{\sqrt{\answer{5}}}\begin{bmatrix}\answer{1} & \answer{-2}\\\answer{2} & \answer{1}\end{bmatrix}$$
\end{problem}

\begin{problem}\label{prob:make_ortho_matrix7}
Normalize the columns to make the following matrix orthogonal.
$A = \begin{bmatrix}
-1 & 2 & 2 \\
2 & -1 & 2 \\
2 & 2 & -1
\end{bmatrix}$
$$\frac{1}{\answer{3}}\begin{bmatrix}\answer{-1} & \answer{2} & \answer{2} \\
\answer{2} & \answer{-1} & \answer{2} \\
\answer{2} & \answer{2} & \answer{-1}
\end{bmatrix}$$
\end{problem}

\begin{problem}\label{prob:findQ}
For each matrix $A$, find an orthogonal matrix $Q$ such that $Q^{-1}AQ$ is diagonal.

\begin{question}
$$A = \begin{bmatrix}
0 & 1 \\
1 & 0
\end{bmatrix}$$
$$Q=\frac{1}{\sqrt{\answer{2}}}\begin{bmatrix}1 & -1\\\answer{1} & \answer{1}\end{bmatrix}$$
\end{question}

\begin{question}
    $$A = \begin{bmatrix}
3 & 0 & 7 \\
0 & 5 & 0 \\
7 & 0 & 3
\end{bmatrix}$$
List eigenvalues in increasing order.
$$\lambda_1=\answer{-4},\quad\lambda_2=\answer{5},\quad\lambda_3=\answer{10}$$
$$Q^{-1}AQ=\begin{bmatrix}\answer{-1}/\sqrt{\answer{2}} & \answer{0} & \answer{1}/\sqrt{\answer{2}}\\\answer{0} & \answer{1} & \answer{0}\\1/\sqrt{\answer{2}} & \answer{0} & 1/\sqrt{\answer{2}}\end{bmatrix}\begin{bmatrix}\lambda_1 & 0 & 0\\0 & \lambda_2 & 0\\0 & 0& \lambda_3\end{bmatrix}\begin{bmatrix}\answer{-1}/\sqrt{\answer{2}} & \answer{0} & \answer{1}/\sqrt{\answer{2}}\\\answer{0} & \answer{1} & \answer{0}\\1/\sqrt{\answer{2}} & \answer{0} & 1/\sqrt{\answer{2}}\end{bmatrix}$$
\end{question}

% \begin{question}
% $$A = \begin{bmatrix}
% 1 & 1 & 0 \\
% 1 & 1 & 0 \\
% 0 & 0 & 2
% \end{bmatrix}$$
% \end{question}

% (challenging problem) $A = \begin{bmatrix}
% 3 & 5 & -1 & 1 \\
% 5 & 3 & 1 & -1 \\
% -1 & 1 & 3 & 5 \\
% 1 & -1 & 5 & 3
% \end{bmatrix}$


%\begin{sol}
%\begin{enumerate} 
%\setcounter{enumi}{1}
%\item  $\frac{1}{\sqrt{2}}\begin{bmatrix}
%1 & -1 \\
%1 & 1
%\end{bmatrix}$

%\setcounter{enumi}{3}
%\item  $\frac{1}{\sqrt{2}}\begin{bmatrix}
%0 & 1 & 1\\
%\sqrt{2} & 0 & 0\\
%0 & 1 & -1
%\end{bmatrix}$

%\setcounter{enumi}{5}
%\item  $\frac{1}{3\sqrt{2}}\begin{bmatrix}
%2\sqrt{2} & 3 & 1\\
%\sqrt{2} & 0 & -4\\
%2\sqrt{2} & -3 & 1
%\end{bmatrix}$ or %$\frac{1}{3}\begin{bmatrix}
%2 & -2 & 1\\
%1 & 2 & 2\\
%2 & 1 & -2
%\end{bmatrix}$

%\setcounter{enumi}{7}
%\item  $\frac{1}{2}\begin{bmatrix}
%1 & -1 & \sqrt{2} & 0\\
%-1 & 1 & \sqrt{2} & 0\\
%-1 & -1 & 0 & \sqrt{2}\\
%1 & 1 & 0 & \sqrt{2}
%\end{bmatrix}$

%\end{enumerate}
%\end{sol}
\end{problem}

\begin{problem}\label{prob:ortho_diag_implies_symmetric}
Suppose $A$ is orthogonally diagonalizable.  Prove that $A$ is symmetric.  (This is the easy direction of the "if and only if" in Theorem \ref{th:PrinAxes}.)
\end{problem}

\begin{problem}\label{ex:8_2_15}
Prove the converse of Theorem~\ref{th:dotpSymmetric}:

If $(A\vec{x}) \dotp \vec{y} = \vec{x} \dotp (A\vec{y})$ for all $n$-columns $\vec{x}$ and $\vec{y}$, then $A$ is symmetric.

%\begin{sol}
%If $\vec{x}$ and $\vec{y}$ are respectively columns $i$ and $j$ of $I_{n}$, then $\vec{x}^{T}A^{T}\vec{y} = \vec{x}^{T}A\vec{y}$ shows that the $(i, j)$-entries of $A^{T}$ and $A$ are equal.
%\end{sol}
\end{problem}



\begin{problem}\label{prob:ortho19}
Let $Q$ be an orthogonal matrix.

\begin{enumerate}
\item Show that $\det Q = 1$ or $\det Q = -1$.

\item Give $2 \times 2$ examples of $Q$ such that $\det Q = 1$ and $\det  Q = -1$.

% \item If $\det  Q = -1$, show that $I + Q$ has no inverse.
% \begin{hint}
% $Q^{T}(I + Q) = (I + Q)^{T}$.
% \end{hint}

% \item If $P$ is $n \times n$ and $\det P \neq (-1)^{n}$, show that $I - P$ has no inverse.

% \begin{hint}
% $P^{T}(I - P) = -(I - P)^{T}$
% \end{hint}
\end{enumerate}
%\begin{sol}
%\begin{enumerate} 
%\setcounter{enumi}{1}
%\item $\det \begin{bmatrix}
%\cos\theta & -\sin\theta \\
%\sin\theta & \cos\theta
%\end{bmatrix} = 1$ \\ and $\det \begin{bmatrix}
%\cos\theta & \sin\theta \\
%\sin\theta & -\cos\theta
%\end{bmatrix} = -1$
%\begin{remark}
%These are the \textit{only} $2 \times 2$ examples.
%\end{remark}

%\setcounter{enumi}{3}
%\item Use the fact that $P^{-1} = P^{T}$ to show that $P^{T}(I - P) = -(I - P)^{T}$. Now take determinants and use the hypothesis that $\det P \neq (-1)^{n}$.

%\end{enumerate}
%\end{sol}
\end{problem}


% \begin{problem}\label{prob:ortho21}
% A matrix that we obtain from the identity matrix by writing its rows in a different order is called a \dfn{permutation matrix} (see Theorem \ref{th:LUPA}). Show that every permutation matrix is orthogonal.
% \end{problem}


% \begin{problem}\label{prob:ortho25}
% Show that the following are equivalent for an $n \times n$ matrix $Q$.


% \begin{enumerate} 
% \item $Q$ is orthogonal.

% \item $\norm{Q\vec{x}} = \norm{\vec{x}}$ for all $\vec{x}\in\RR^n$.

% \item $\norm{ Q\vec{x} - Q\vec{y}} = \norm{\vec{x} - \vec{y}}$ for all $\vec{x}$, $\vec{y}\in \RR^n$.

% \item $(Q\vec{x}) \dotp (Q\vec{y}) = \vec{x} \dotp \vec{y}$ for all columns $\vec{x}$, $\vec{y}\in\RR^n$.

% \begin{hint}
% %For (c) $\Rightarrow$ (d), see Exercise \ref{ex:5_3_14}(a). 
% For (d) $\Rightarrow$ (a), show that column $i$ of $Q$ equals $Q\vec{e}_{i}$, where $\vec{e}_{i}$ is column $i$ of the identity matrix.
% \end{hint}
% \end{enumerate}

% \begin{remark}
%     This exercise shows that linear transformations with orthogonal standard matrices are distance-preserving (b,c) and angle-preserving (d).
% \end{remark}

% \end{problem}




% \begin{problem}\label{prob:SchurChallenge}
% Modify the proof of Theorem~\ref{th:PrinAxes} to prove Theorem \ref{th:Schur}.
%If $A\vec{x}_{1} = \lambda_{1}\vec{x}_{1}$ where $\norm{\vec{x}_{1}} = 1$, let $\{\vec{x}_{1}, \vec{x}_{2}, \dots, \vec{x}_{n}\}$ be an orthonormal basis of $\RR^n$, and let $P_{1} = \begin{bmatrix}
%\vec{x}_{1} & \vec{x}_{2} & \cdots &  \vec{x}_{n}
%\end{bmatrix}$. Then $P_{1}$ is orthogonal and $P_{1}^TAP_{1} = \begin{bmatrix}
%\lambda_{1} & B \\
%0 & A_{1}
%\end{bmatrix}$ in block form. By induction, let $Q^{T}A_{1}Q = T_{1}$ be upper triangular where $Q$ is of size $(n-1)\times(n-1)$ and orthogonal. Then $P_{2} = \begin{bmatrix}
%1 & 0 \\
%0 & Q
%\end{bmatrix}$ is orthogonal, so $P = P_{1}P_{2}$ is also orthogonal and $P^TAP = \begin{bmatrix}
%\lambda_{1} & BQ \\
%0 & T_{1}
%\end{bmatrix}$
% is upper triangular.
% \end{problem}

\section*{Text Source} This section was adapted from Section 8.2 of Keith Nicholson's \href{https://open.umn.edu/opentextbooks/textbooks/linear-algebra-with-applications}{\it Linear Algebra with Applications}. (CC-BY-NC-SA)

W. Keith Nicholson, {\it Linear Algebra with Applications}, Lyryx 2018, Open Edition, p. 424


\end{document}
